\documentclass[10pt, letterpaper]{article}

\usepackage{../../../../template/template}

% ------------------------
% Impostazioni documento
% ------------------------

\setDocTitle{Verbale esterno 26/03/2026}
\setDocVersion{-}
\setDocData{26/03/2026}
\setDocIndex{ve\_2026\_03\_26}
\setDocState{2}
\setDocRecipient{2}

\begin{document}

\makeTitlePage

\newpage

\tableofcontents

\newpage

% -------------------------------------------------
% Informazioni riunione
% -------------------------------------------------
\makeMeetingInfoTable{26/03/2026}{16:30}{17:00}{via \textit{Microsoft Teams$_G$}}

% ------------------------
% Partecipanti
% ------------------------
\addParticipant{Filippo}{Venzo}{Programmatore}{P}
\addParticipant{Valeria}{Baleanu}{Programmatore}{P}
\addParticipant{Luca}{Granziero}{Amministratore}{P}
\addParticipant{Christian}{Libralato}{Responsabile}{P}
\addParticipant{Francesco}{Pasqual}{Programmatore}{P}
\addParticipant{Leonardo}{Pellizzon}{Programmatore}{P}
\addParticipant{Giuseppe}{De Fina}{Verificatore}{A}
\makeMeetingParticipantsTable

% -------------------------------------------------
% Ordine del giorno
% -------------------------------------------------
\section{Ordine del giorno}
\begin{itemize}
      \item Avanzamento del Progetto e Scadenze;
      \item implementazione Suggerimenti e Analytics;
      \item gestione della Cache degli Impianti;
      \item allineamento Frontend e Backend.
\end{itemize}

% -------------------------------------------------
% Approfondimento
% -------------------------------------------------
\section{Approfondimento}

\subsection*{Avanzamento del Progetto e Scadenze}
Il Responsabile ha illustrato lo stato dei lavori, riportando che buona parte della logica di \textit{backend} e \textit{frontend} è stata implementata, ad esclusione
della parte di analytics, notifiche e definizione estetica dell'applicazione. 

Riguardo la documentazione è stato annunciato l'inizio della stesura della
Specifica tecnica e l'ultimazione dei diagrammi UML e C4, a meno di ultime rifiniture.
La Proponente ha infine esposto la necessità di fornire un data per fissare il collaudo dell'applicativo, che coinvolgerà anche il team
marketing, da comunicare entro fine settimana.

\subsection*{Implementazione Suggerimenti e Analytics}
Il team ha chiesto consigli su come implementare i suggerimenti, sono state valutate delle strategie basate su  approcci statici, di machine learning e con l'utilizzo di API di LLM,
suggerendo di partire da suggerimenti statici e, in parallelo, sperimentare soluzioni più dinamiche.

In particolare, in riferimento all'utilizzo di API per LLM, è stato spiegato che questa soluzione, che si basa sul passaggio di dati di riferimento e dati da analizzare tramite prompt engineering, può essere rapida e flessibile, mentre l'addestramento di un modello ML classico richiederebbe più tempo e dati.

È stato successivamente chiarito un dubbio legato alla soglia da utilizzare per considerare una presenza "prolungata". L'esito è stato che la soglia può variare in base al contesto ma può corrispondere a valori al di sotto di un'ora per situazioni più critiche, come la presenza in bagno.

È infine stato offerto supporto per l'accesso a modelli come Amazon Nova qualora necessario, suggerendo ad ogni modo di provare con i modelli conosciuti.


\subsection*{Gestione della Cache degli Impianti}

È stato illustrato il funzionamento attuale della cache degli impianti e sono state fornite indicazioni su come ottimizzare l'aggiornamento tramite subscription alle notifiche di modifica, evitando l'aggiornamento periodico e riducendo l'overhead. Sono dunque state introdotte le subscription di tipo 'node',
che notificano il cambiamento della struttura dell'impianto, permettendo di invalidare e aggiornare la cache in modo mirato.

\subsection*{Allineamento Frontend e Backend}
Il team ha infine presentato una demo del sistema funzionante, frutto di una prima integrazione tra \textit{backend} e \textit{frontend}.
È stato confermato che la documentazione delle API tramite file OpenAPI (swagger.yaml) è sufficiente e che la versione 3.x è accettata. 

Inoltre, riguardo allo stile grafico è stato suggerito di usare librerie di icone e uno stile moderno, in riferimento al quale verranno a breve 
condivisi i commenti elaborati dal team UX/UI sui bozzetti grafici realizzati.



% -------------------------------------------------

% Decisioni
% -------------------------------------------------
\addDecision{I suggerimenti saranno realizzati in maniera ibrida con l'utlizzo di metodi statici e di API per LLM.}
\makeDecisionTable

% -------------------------------------------------
% Azioni
% -------------------------------------------------
\addTodo{\noOne}{Comunicare entro il fine settimana una data per il collaudo.}{C. Libralato}{27/03/2026}

\makeTodoTable

\end{document}
